\documentclass[11pt,a4paper]{article}
\usepackage[margin=2.5cm]{geometry}
\usepackage[T1]{fontenc}
\usepackage{lmodern}
\usepackage{amsmath,amssymb,amsthm}
\usepackage{booktabs}
\usepackage{graphicx}
\usepackage[colorlinks=true,linkcolor=blue,citecolor=blue,urlcolor=blue]{hyperref}
\usepackage{xcolor}
\newtheorem{result}{Result}
\newtheorem{corollary}{Corollary}
\theoremstyle{remark}
\newtheorem{remark}{Remark}
\numberwithin{equation}{section}

\title{Universal small-spin formulas for the shadow area, the capture cross-section and the spin-down of Kerr-like rotating black holes}
\author{Jacob Goodchild\thanks{Independent researcher. Contact: jacob.l.goodchild@gmail.com.
This work was carried out with substantial assistance from an AI system; see the
AI disclosure statement at the end of the paper.}}
\date{October 2026}

\begin{document}
\maketitle

\begin{abstract}
Many rotating black-hole metrics studied in the literature are generated from a static,
spherically symmetric seed $ds^2=-f\,dt^2+dr^2/f+r^2d\Omega^2$ by the Newman--Janis algorithm in the
form of Azreg-A\"{\i}nou, which yields the Kerr metric with $2Mr$ replaced by $r^2(1-f(r))$. For this
whole class we give closed formulas, in terms of $f$ and its derivatives at a single radius, for
the leading effects of a small spin $a$ on three observables: (i) the area of the black-hole
shadow seen by a distant observer at any inclination, to order $a^2$ and, with a term-by-term
finite-part method, to orders $a^4$, $a^6$ and $a^8$; (ii) the capture cross-section for neutral
massive particles of any speed arriving from any direction, to order $a^2$ (and to order $a^4$
along the spin axis); and (iii) the mean angular momentum swallowed per unit captured energy, which
controls the spin-down of a black hole accreting collisionless matter. The order-$a^2$ shadow formula
depends on $f$ only through $f(R)$ and the curvature combination $\kappa=2f-R^2f''$ at the photon
sphere $R$; the massive-particle formula contains the combination $3ff'+rff''-2rf'^2$ whose zero
defines the innermost stable circular orbit. Specialising to Kerr--Newman gives explicit expressions valid at arbitrary charge, including the shadow area
to order $a^4$; for Kerr the method reproduces the known low-order terms and yields the order-$a^6$ and $a^8$
coefficients of the shadow area at any inclination. All results
were checked against direct high-precision computations of the exact shadow contour or of the
exact capture boundary for Kerr, Kerr--Newman, rotating Bardeen and rotating Hayward black holes,
with agreement to 11--22 significant digits.
\end{abstract}

\section{Introduction}

The shadow of a black hole---the dark region in the sky bounded by the image of the unstable
photon orbits---is now an observable~\cite{EHT2019}. For the Kerr metric the boundary of the shadow
is known in closed parametric form~\cite{Bardeen1973,Chandrasekhar1983}. A large body of work
constructs rotating counterparts of non-Kerr static black holes (regular black holes, black holes
surrounded by matter, quantum-corrected black holes, and so on) with the Newman--Janis
algorithm~\cite{NewmanJanis1965}, most often in the non-complexified form of
Azreg-A\"{\i}nou~\cite{AzregAinou2014}, and then computes their shadows. For this class the
geodesic equations separate and general formulas for the shadow contour are
available~\cite{Tsukamoto2018,Shaikh2019}; observables such as the area, the
circumference or the displacement of the shadow are then usually computed numerically, model by
model.

The purpose of this paper is to give \emph{universal} closed formulas for the small-spin behaviour
of such observables, valid for any seed function $f(r)$. The main results are:
\begin{itemize}
\item the order-$a^2$ correction to the shadow area at any inclination, eq.~\eqref{eq:A2};
\item the order-$a^4$ correction, eqs.~\eqref{eq:A4}--\eqref{eq:N2}, obtained with a
finite-part (residue at infinity) method that extends mechanically to $a^6$ and $a^8$;
\item the order-$a^2$ correction to the capture cross-section of neutral massive particles of
arbitrary speed and direction, eq.~\eqref{eq:sig2}, and its on-axis continuation to order $a^4$,
eq.~\eqref{eq:sigpole};
\item the mean angular momentum per unit energy carried into the hole by captured particles,
eq.~\eqref{eq:torque}.
\end{itemize}
For Kerr and Kerr--Newman these specialise to explicit rational functions of the radius of the
relevant circular orbit (Section~\ref{sec:special}). Several of these special cases were known before, in
particular for Kerr~\cite{Bozza2006,GrallaLupsasca2020shape,KobialkoGaltsov2026,MachMomenniaSarbach2026}; Section~\ref{sec:prior}
lists precisely which results are previously known and which, to the best of our knowledge, are new.

Units are $G=c=1$. Lengths are measured in units of the mass parameter of the seed metric where
relevant; all formulas below are invariant under the rescaling $r\to\lambda r$, $a\to\lambda a$,
area $\to\lambda^2$ area, with $f(r)$ rescaled accordingly.

\section{Setting}\label{sec:setting}

\subsection{Metric and geodesics}
Let $f(r)$ be a smooth function and put
\begin{equation}
\Sigma=r^2+a^2\cos^2\vartheta,\qquad \Delta=r^2f(r)+a^2 .
\end{equation}
The rotating metrics considered here are
\begin{multline}
ds^2=-\Big(1-\frac{r^2(1-f)}{\Sigma}\Big)dt^2-\frac{2a r^2(1-f)\sin^2\vartheta}{\Sigma}\,dt\,d\varphi
+\frac{\Sigma}{\Delta}dr^2+\Sigma\, d\vartheta^2\\
+\frac{\big[(r^2+a^2)^2-\Delta a^2\sin^2\vartheta\big]\sin^2\vartheta}{\Sigma}d\varphi^2 ,
\label{eq:metric}
\end{multline}
which for $f=1-2M/r$ is the Kerr metric and for $f=1-2M/r+Q^2/r^2$ the Kerr--Newman metric. This is
the form produced by the Newman--Janis algorithm without complexification~\cite{AzregAinou2014}
when the seed is $ds^2=-f\,dt^2+dr^2/f+r^2d\Omega^2$. The Hamilton--Jacobi equation separates, and for
a geodesic with energy $E$, axial angular momentum $L$, Carter constant $Q$ and mass $\mu$
(all per unit mass for $\mu=1$; $\mu=0$ for light),
\begin{align}
\Sigma\,\frac{dr}{d\lambda}&=\pm\sqrt{\mathcal R(r)},&
\mathcal R(r)&=\big[E(r^2+a^2)-aL\big]^2-\Delta\big[\mu^2r^2+(L-aE)^2+Q\big],\label{eq:R}\\
\Sigma\,\frac{d\vartheta}{d\lambda}&=\pm\sqrt{\Theta(\vartheta)},&
\Theta(\vartheta)&=Q-\cos^2\vartheta\Big[a^2(\mu^2-E^2)+\frac{L^2}{\sin^2\vartheta}\Big].
\end{align}

\subsection{The impact plane}
Consider particles that arrive from infinity (or reach a distant observer) along the direction
making angle $\vartheta_o$ with the spin axis, and write
\begin{equation}
C\equiv\cos^2\vartheta_o .
\end{equation}
With $p=\sqrt{E^2-\mu^2}$ the asymptotic momentum, the Cartesian impact parameters (celestial
coordinates for light) are $\alpha=-L/(p\sin\vartheta_o)$, $\beta=\pm\sqrt{\Theta(\vartheta_o)}/p$, so that
\begin{equation}
b^2\equiv\alpha^2+\beta^2=\frac{L^2+Q}{p^2}+a^2C .
\label{eq:b2}
\end{equation}
The \emph{shadow} (for light) and the \emph{capture region} (for massive particles of given $E$)
are the regions of the $(\alpha,\beta)$ plane whose geodesics reach the horizon. Their boundaries
are traced by orbits on which $\mathcal R$ has a double root, $\mathcal R=\mathcal R'=0$; we write
$\mathcal A$ for the area of the shadow and $\sigma$ for the capture cross-section.

\subsection{The non-rotating data}
At $a=0$ the boundary is a circle. For light it is set by the photon sphere $r=R$,
\begin{equation}
Rf'(R)=2f(R),\qquad \mathcal A_0=\frac{\pi R^2}{f(R)} .
\label{eq:photon}
\end{equation}
For massive particles with specific energy $E>1$ (speed $v$ at infinity, $E^2=1/(1-v^2)$) it is
set by the unstable circular orbit of that energy, of radius $r$, with
\begin{equation}
E^2=\frac{2f^2}{2f-rf'},\qquad L^2=\frac{r^3f'}{2f-rf'},\qquad
\sigma_0=\frac{\pi L^2}{E^2-1}=\frac{\pi r^3 f'}{2f^2-2f+rf'} ,
\label{eq:circ}
\end{equation}
where $f$ and its derivatives are evaluated at $r$. These relations are standard; for the
Schwarzschild metric they give the classical capture cross-section~\cite{Unruh1976} $\sigma_0=\pi r^3/(4-r)$ (with
$M=1$), which tends to $16\pi/v^2$ for slow particles and to $27\pi$ for light. Throughout, $r$
denotes this circular-orbit radius, which lies between the photon sphere (light) and the marginally
bound orbit ($v\to0$).

We use the abbreviations
\begin{equation}
\kappa\equiv2f-R^2f''(R)\quad\text{(at the photon sphere)},\qquad
J\equiv3ff'+rff''-2rf'^2\quad\text{(at radius $r$)} .
\label{eq:kappaJ}
\end{equation}
For static metrics, $\kappa$ is related to the Lyapunov exponent $\lambda$ of the photon sphere
(in coordinate time) by $\lambda^2=f\kappa/(2R^2)$~\cite{Cardoso2009}, and $J=0$ is the condition for
the innermost stable circular orbit. For Schwarzschild, $\kappa=2$ and $J=2(r-6)/r^3$.

\section{Method}\label{sec:method}

\subsection{Perturbation of the double-root conditions}
For the order-$a^2$ results we parametrise the boundary orbits by their total angular momentum
$\ell=\sqrt{L^2+Q}$ and the inclination $\cos i=L/\ell$, and expand the critical values
\begin{equation}
\ell_c=\ell_0+a\,\ell_1\cos i+a^2\ell_2(\cos^2 i)+O(a^3),\qquad r_c=r+O(a),
\end{equation}
by solving $\mathcal R=\mathcal R'=0$ order by order, with $f$ replaced by its Taylor series about the
unperturbed radius. To leading order the inclination is related to the polar angle $\phi$ in the
impact plane by $\cos i=-\sin\vartheta_o\cos\phi$, with $\phi$ measured from the $\alpha$ axis. The area of a star-shaped region is
$\frac12\oint b^2\,d\phi$, and by \eqref{eq:b2}, $b^2=\ell_c^2/p^2+a^2C$. Averaging over $\phi$ gives
the order-$a^2$ coefficient; the $O(a^2)$ reparametrisation between $\phi$ and $i$ does not contribute
at this order because the $O(a)$ term is odd in $\cos\phi$. In particular one finds the universal
first-order coefficient
\begin{equation}
\ell_1=-E\,\frac{1-f}{f}
\label{eq:ell1}
\end{equation}
(sign chosen so that prograde orbits have $\cos i>0$). The derivative $f'''$ drops out of all
order-$a^2$ results. The computations were done with \textsc{SymPy}.

\subsection{Finite-part expansion of the shadow area}\label{sec:fp}
For light the boundary of the shadow is known in closed parametric form for any
$\Delta(r)$~\cite{Tsukamoto2018}:
\begin{equation}
\xi(r)=\frac{(r^2+a^2)\Delta'-4r\Delta}{a\Delta'},\qquad
\eta(r)=\frac{r^2\big[16a^2\Delta-(4\Delta-r\Delta')^2\big]}{a^2\Delta'^2},
\label{eq:xieta}
\end{equation}
with $\xi=L/E$, $\eta=Q/E^2$, $\alpha=-\xi/\sin\vartheta_o$ and
$\beta^2=\eta+a^2C-\xi^2\cot^2\vartheta_o$. Set $r=R+a x$. Because $4\Delta-r\Delta'=r^2(2f-rf')+4a^2$
vanishes at $a=0$, $r=R$, every coefficient of the expansion of $\xi$ and $\eta$ in powers of $a$ at
fixed $x$ is a \emph{polynomial} in $x$. After the linear change of variable $u=\alpha_0(x)$
(the zeroth-order abscissa), $\beta^2=b_0^2-u^2+\delta(u,a)$ with $b_0^2=R^2/f$ and $\delta$ a
polynomial in $u$ at every order in $a$. The area
$\mathcal A=2\int\beta\,d\alpha$ is an analytic function of $a$, and its expansion can be computed term by term
with the finite-part (Hadamard) rule
\begin{equation}
\mathrm{FP}\!\int_{-b_0}^{b_0}u^m\,(b_0^2-u^2)^{\frac12-j}\,du
=-\pi\,(-1)^j\binom{\frac12-j}{n}(-b_0^2)^n,\qquad n=\frac{m+2-2j}{2},
\label{eq:FP}
\end{equation}
for even $m$ with $n\ge0$ (and zero otherwise), which follows from the residue of
$u^{m}(b_0^2-u^2)^{1/2-j}$ at infinity. Expanding $\sqrt{b_0^2-u^2+\delta}\;d\alpha/du$ binomially in
$\delta$ and applying \eqref{eq:FP} to each monomial gives the area to any order. With truncated power
series over the rational function field in $f(R),R f'(R),\dots$ this takes about one second at
order $a^4$, ten seconds at order $a^6$ and seven minutes at order $a^8$ on a laptop-class CPU.

\section{Results}\label{sec:results}

Throughout this section $f,f',f'',\dots$ denote values at the photon sphere $R$ (for the shadow) or at
the circular-orbit radius $r$ (for massive particles), and $C=\cos^2\vartheta_o$.

\subsection{Shadow area to order \texorpdfstring{$a^2$}{a2}}
\begin{result}[Shadow area]\label{res:A2}
\begin{equation}
\mathcal A=\frac{\pi R^2}{f}\left[1+a^2\left\{(1-C)\left(\frac{1}{2fR^2}-\frac{4}{R^2\kappa}\right)
+C\,\frac{2f-1}{fR^2}\right\}+O(a^4)\right].
\label{eq:A2}
\end{equation}
\end{result}
\begin{corollary}[Polar observer]
For $\vartheta_o=0$ the shadow is a disc, of area
$\mathcal A=\pi\big[R^2/f+a^2(2f-1)/f^2\big]+O(a^4)$. (The exact polar radius of integrable
spinning spacetimes is given in~\cite{Salehi2024}; this is its small-spin expansion.)
\end{corollary}
\begin{corollary}[Displacement and circumference]
To first order in $a$ the shadow is displaced along the $\alpha$ axis by $a\sin\vartheta_o\,(1-f)/f$
(for Kerr: $2a\sin\vartheta_o$, a known result~\cite{Bozza2006,GrallaLupsasca2020shape}). To order $a^2$ the boundary is a displaced circle plus a
$\cos2\phi$ deformation, which changes neither area nor length at this order; hence the
circumference is $P=2\pi R f^{-1/2}\big[1+\tfrac12 a^2\{\cdots\}+O(a^4)\big]$ with the same bracket
$\{\cdots\}$ as in \eqref{eq:A2}.
\end{corollary}

\subsection{Shadow area to order \texorpdfstring{$a^4$}{a4}}
Write $f_3=R^3f'''(R)$ and $f_4=R^4f''''(R)$.
\begin{result}[Next order]\label{res:A4}
\begin{equation}
\mathcal A=\frac{\pi R^2}{f}\Big[1+a^2K_2+\frac{a^4}{R^4}\,\frac{N_0+N_1C+N_2C^2}{8f^2\kappa^5}+O(a^6)\Big],
\label{eq:A4}
\end{equation}
where $K_2$ is the bracket of \eqref{eq:A2} and
\begin{align}
N_0&=-\kappa^5+16f\kappa^4-160f^2\kappa^3+32f^2f_3\kappa^2-384f^3\kappa^2+256f^3f_3\kappa-32f^3f_4\kappa-96f^3f_3^2,\label{eq:N0}\\
N_1&=2\big[-3\kappa^5+4f\kappa^5+32f\kappa^4-32f^2\kappa^4-160f^2\kappa^3+32f^2f_3\kappa^2+192f^3\kappa^3\nonumber\\
&\qquad+384f^3\kappa^2-64f^3f_3\kappa^2-256f^3f_3\kappa+32f^3f_4\kappa+96f^3f_3^2\big],\label{eq:N1}\\
N_2&=15\kappa^5-24f\kappa^5-144f\kappa^4+8f^2\kappa^5+192f^2\kappa^4+480f^2\kappa^3-96f^2f_3\kappa^2\nonumber\\
&\qquad-64f^3\kappa^4-384f^3\kappa^3-384f^3\kappa^2+128f^3f_3\kappa^2+256f^3f_3\kappa-32f^3f_4\kappa-96f^3f_3^2.\label{eq:N2}
\end{align}
\end{result}
The coefficients of $a^6$ and $a^8$ are obtained by the same procedure; they are long rational
expressions in $f$, $R^kf^{(k)}(R)$ ($k\le6$ and $k\le8$) and $C$, and are provided in machine-readable
form with the accompanying code (see Code and data availability).

\subsection{Capture cross-section to order \texorpdfstring{$a^2$}{a2}, any speed}
\begin{result}[Capture]\label{res:sig2}
For neutral particles with speed $v$ at infinity, arriving at angle $\vartheta_o$ to the spin axis,
with $r$ the radius of the unstable circular orbit of energy $E=(1-v^2)^{-1/2}$ in the seed metric
[eq.~\eqref{eq:circ}],
\begin{equation}
\sigma=\sigma_0\left[1+a^2\left\{(1-C)\left(\frac{1}{2fr^2}+\frac{f'^2}{rfJ}\right)
+C\,\frac{2f+rf'-2}{r^3f'}\right\}+O(a^4)\right].
\label{eq:sig2}
\end{equation}
\end{result}
At the photon sphere ($rf'=2f$, hence $J=-f\kappa/r$) the bracket reduces to that of
\eqref{eq:A2}, as it must.

\begin{result}[On the spin axis to order $a^4$]\label{res:pole}
For $\vartheta_o=0$ the capture region is exactly a disc of radius$^2$ $Q/p^2+a^2$ (with $L=0$), and
\begin{equation}
\sigma=\sigma_0\left[1+a^2\,\frac{2f+rf'-2}{r^3f'}+a^4\,\frac{2(1-f)^2(3f'+rf'')}{r^5f'J}+O(a^6)\right].
\label{eq:sigpole}
\end{equation}
\end{result}

\subsection{Angular momentum swallowed}
Let $\langle L\rangle$ be the mean axial angular momentum of the captured particles (averaged with
uniform weight over the capture region of the impact plane) and $E$ their energy, both per unit
rest mass. Then $\langle L\rangle/E=dJ_{\rm BH}/dM_{\rm BH}$ is the angular momentum gained per unit
mass-energy gained.
\begin{result}[Spin-down]\label{res:torque}
\begin{equation}
\frac{\langle L\rangle}{E}=-a\sin^2\vartheta_o\,\frac{1-f}{f}+O(a^3),
\label{eq:torque}
\end{equation}
for massive particles of any speed ($f$ at $r$) and for light ($f$ at $R$).
\end{result}
For a black hole immersed in an isotropic flux of particles of a single speed, averaging over
directions with the leading-order weight $\sigma_0$ gives $\langle L\rangle/E=-\tfrac23a(1-f)/f$ and,
for the dimensionless spin $a_*=J_{\rm BH}/M_{\rm BH}^2$,
\begin{equation}
\frac{da_*}{d\ln M_{\rm BH}}=-2a_*-\frac{2a_*}{3}\,\frac{1-f}{f}+O(a_*^3).
\end{equation}

\section{Special cases}\label{sec:special}

\subsection{Kerr}
With $f=1-2/r$ ($M=1$): $R=3$, $f(R)=1/3$, $\kappa=2$, $f_3=4$, $f_4=-16$. The particle
speed is encoded by the circular-orbit radius
\begin{equation}
r=\frac{4v^2-1+\sqrt{1+8v^2}}{2v^2}\in[3,4],\qquad \sigma_0=\frac{\pi r^3}{4-r}.
\end{equation}
Results~\ref{res:A2}--\ref{res:torque} give the following. The $a^2$ and $a^4$ terms of \eqref{eq:kerrA} are
known~\cite{Bozza2006,GrallaLupsasca2020shape,KobialkoGaltsov2026}; the $a^6$ and $a^8$ terms do not appear in the literature we found. The
capture formula \eqref{eq:kerrsig} at order $a^2$ and the order-$a$ term of $\langle L\rangle/E$ follow from the slow-rotation
expansion of the critical Carter constant of Mach, Momennia and Sarbach~\cite{MachMomenniaSarbach2026} (who display
the direction-averaged versions), and agree with the fitted coefficients of~\cite{Karydas2026}:
\begin{align}
\frac{\mathcal A}{27\pi}&=1-\frac{(1+C)a^2}{18}-\Big(\frac1{72}+\frac{5C}{972}-\frac{5C^2}{1944}\Big)a^4
-\Big(\frac{2059}{314928}+\frac{167C}{104976}-\frac{185C^2}{104976}-\frac{5C^3}{3888}\Big)a^6\nonumber\\
&\quad-\Big(\frac{88175}{22674816}+\frac{3869C}{5668704}-\frac{4865C^2}{3779136}-\frac{83C^3}{69984}-\frac{667C^4}{7558272}\Big)a^8+O(a^{10}),\label{eq:kerrA}\\
\frac{\sigma}{\sigma_0}&=1-a^2\,\frac{r+3C(4-r)}{2r^2(6-r)}+O(a^4),\qquad
\frac{\sigma(\vartheta_o=0)}{\sigma_0}=1-\frac{a^2}{r^2}-\frac{4a^4}{r^4(6-r)}+O(a^6),\label{eq:kerrsig}\\
\frac{\langle L\rangle}{E}&=-\frac{2a\sin^2\vartheta_o}{r-2}+O(a^3).
\end{align}
For light ($r=3$) the capture result reproduces the $a^2$ term of \eqref{eq:kerrA}; for slow particles
($r\to4$) the $a^2$ correction $-a^2/16$ is independent of direction.

\paragraph{Higher orders for Kerr.} By fitting high-precision
evaluations of the exact Kerr capture region at radii $r\in\{15/4,7/2,10/3,13/4,18/5,16/5\}$ we
identified exact rational coefficients for the full order-$a^4$ term and for the order-$a^3$ term of
the swallowed angular momentum:
\begin{align}
\frac{\sigma}{\sigma_0}&=1-a^2\frac{r+3C(4-r)}{2r^2(6-r)}
+a^4\,\frac{\alpha_8+\beta_8C+\gamma_8C^2}{8r^4(6-r)^5}+O(a^6),\label{eq:kerr4}\\
\alpha_8&=-r^2(2r^3+42r^2-549r+1458),\quad
\beta_8=6r(2r^4+2r^3-239r^2+1302r-2016),\nonumber\\
\gamma_8&=-10r^5-2r^4+1653r^3-13266r^2+39744r-41472,\nonumber\\
\frac{\langle L\rangle}{E}&=-\frac{2a(1-C)}{r-2}
+a^3(1-C)\,\frac{\tfrac32(r-5)-C\,\frac{7r^2-63r+144}{2r}}{(r-2)(6-r)^3}+O(a^5).\label{eq:kerrL3}
\end{align}
Eq.~\eqref{eq:kerr4} satisfies $\alpha_8+\beta_8+\gamma_8=-32(6-r)^4$, reproducing the on-axis
result~\eqref{eq:kerrsig}, which was derived analytically; at $r=3$ it reproduces the $a^4$ term
of \eqref{eq:kerrA}. Eq.~\eqref{eq:kerrL3} was found from data at $C=0$ and $C=1/2$ and then
confirmed at $C=1/4$ and $C=4/5$; both formulas were further confirmed at radii and angles not used in
their identification ($r=3.55$, $40^\circ$ and $r=3.15$, $70^\circ$; Table~\ref{tab:checks}). These two formulas are identifications
supported by strong numerical evidence, not derivations. The order-$a^4$ capture term is beyond the
expansion of~\cite{MachMomenniaSarbach2026}, which is given to third order in the spin, and has only been available as numerical fits~\cite{Karydas2026};
the $a^3$ term of $\langle L\rangle/E$ can in principle be derived from the third-order terms of~\cite{MachMomenniaSarbach2026}, where it is not displayed. Averaging \eqref{eq:kerrL3} over an
isotropic flux, with the order-$a^2$ weight \eqref{eq:kerrsig}, gives
\begin{equation}
\frac{\langle L\rangle}{E}\Big|_{\rm iso}=-\frac{4a}{3(r-2)}
+\frac{4a^3(3r^3-19r^2+48r-144)}{15r^2(r-2)(6-r)^3}+O(a^5),
\end{equation}
equal to $-(2a/3+a^3/15)$ for slow particles, as found in~\cite{MachMomenniaSarbach2026,MachMomenniaSarbach2026b},
and $-(4a/3+8a^3/81)$ for light.

\subsection{Kerr--Newman}
With $f=1-2/r+q^2/r^2$ ($M=1$, charge $q$; neutral test particles):
\begin{align}
\frac{\sigma}{\sigma_0}&=1-a^2\,\frac{r^2(r-q^2)-C\,(3r^3-12r^2-q^2r^2+18q^2r-8q^4)}
{2r(r-q^2)(6r^2-r^3-9q^2r+4q^4)}+O(a^4),\label{eq:KNsig}\\
\sigma_0&=\frac{\pi r^4(r-q^2)}{4r^2-r^3-4q^2r+q^4},\nonumber\\
\frac{\langle L\rangle}{E}&=-a\sin^2\vartheta_o\,\frac{2r-q^2}{r^2-2r+q^2}+O(a^3),
\end{align}
with $E^2=(r^2-2r+q^2)^2/[r^2(r^2-3r+2q^2)]$. The factor $r^3-6r^2+9q^2r-4q^4$ is the
Reissner--Nordstr\"om ISCO polynomial. For the shadow, with $R=\tfrac12\big(3+\sqrt{9-8q^2}\big)$,
\begin{equation}
\mathcal A=\frac{\pi R^4}{R^2-2R+q^2}\Big[1-a^2\frac{R+3C(R-2)}{R^2(R-1)(2R-3)}+a^4K_4+O(a^6)\Big],
\label{eq:KNA}
\end{equation}
where
\begin{align*}
K_4=\frac{1}{2R^4(R-1)^2(2R-3)^5}\Big[&C^2(R^6-38R^5+286R^4-768R^3+774R^2-108R-162)\\
&-C(18R^6-108R^5+288R^4-492R^3+504R^2-216R)\\
&-15R^6+58R^5-78R^4+36R^3\Big].
\end{align*}
Expanding in $q$ reproduces the double series
$\mathcal A/(27\pi)=1-q^2/3-(1+C)a^2/18-q^4/27-(3+C)a^2q^2/81+\dots$, which is contained in the perturbative
results of Kobialko and Gal'tsov~\cite{KobialkoGaltsov2026}; eqs.~\eqref{eq:KNsig}--\eqref{eq:KNA} are exact in $q$.

\subsection{Regular black holes}
For the rotating Bardeen ($f=1-2r^2/(r^2+g^2)^{3/2}$~\cite{AyonBeatoGarcia2000}) and Hayward ($f=1-2r^2/(r^3+2l^2)$~\cite{Hayward2006}) black holes,
the order-$a^2$ coefficient of $\mathcal A/\mathcal A_0$ at $\vartheta_o=90^\circ$ is $-0.0797470530672$
($g=0.4$) and $-0.0863835115987$ ($l=0.5$) respectively; further values are listed in
Table~\ref{tab:checks}.

\section{Verification}\label{sec:checks}

Every result was checked against an independent computation that does not use the perturbative
expansions.
\begin{itemize}
\item \emph{Shadow:} the exact contour \eqref{eq:xieta} was integrated by tanh--sinh quadrature at 40
digits for $a=0.01,\dots,0.05$, and the coefficients of $a^2$ and $a^4$ were extracted by polynomial
fitting in $a^2$.
\item \emph{Capture:} the exact boundary $b(\phi)$ was computed at 48 impact-plane angles by solving
$\mathcal R=\mathcal R'=0$ with Newton's method at 40 digits, the area evaluated as
$\frac12\oint b^2d\phi$ (trapezoidal rule, spectrally accurate for periodic integrands), and the
coefficients extracted as above. The swallowed angular momentum was computed as the first moment
of the same region.
\item \emph{Symbolic:} the Kerr--Newman specialisation of \eqref{eq:sig2} was checked to agree identically
with an independent Kerr--Newman calculation; the light limit of \eqref{eq:sigpole} agrees
with \eqref{eq:A4} at $C=1$ (exactly for Kerr, where both give $-4a^4/243$, and to $10^{-25}$ for Bardeen); the Kerr limits of the $a^6$ and $a^8$ shadow coefficients agree exactly with
rational coefficients identified independently from 250-digit evaluations of an exact
elliptic-integral formula for the edge-on Kerr shadow area and from a single-integral representation
at general inclination.
\end{itemize}

\begin{table}[ht]
\centering\footnotesize
\resizebox{\textwidth}{!}{%
\begin{tabular}{llll}
\toprule
Quantity & Case & Direct computation & Formula\\
\midrule
$a^2$ coeff.\ of $\mathcal A/\mathcal A_0$ & Bardeen $g=0.4$, $90^\circ$ & $-0.07974705306715083$ & $-0.07974705306715083$\\
 & Bardeen $g=0.4$, $50^\circ$ & $-0.1035696370499696$ & $-0.1035696370499696$\\
 & Hayward $l=0.5$, $90^\circ$ & $-0.08638351159869441$ & $-0.08638351159869439$\\
 & Hayward $l=0.5$, $50^\circ$ & $-0.10712747758$ & $-0.107127477578$\\
 & Kerr, $50^\circ$ & $-0.07850977284258527$ & $-0.07850977284258527$\\
$a^4$ coeff.\ of $\mathcal A/\mathcal A_0$ & Bardeen $g=0.4$, $90^\circ$ & $-0.0317475715025$ & $-0.0317475715025$\\
 & Bardeen $g=0.4$, $50^\circ$ & $-0.0319437641728$ & $-0.0319437641729$\\
 & Kerr--Newman $q=0.6$, $60^\circ$ & $-0.0344681711527$ & $-0.0344681711528$\\
 & Kerr--Newman $q=0.8$, $90^\circ$ & $-0.0918348939951$ & $-0.0918348939966$\\
$a^2$ coeff.\ of $\sigma/\sigma_0$ & Bardeen, $r=3.4$, $60^\circ$ & $-0.08155531033444$ & $-0.08155531033444$\\
 & Bardeen, $r=3.1$, $90^\circ$ & $-0.07771672446027$ & $-0.07771672446026$\\
 & Hayward, $r=3.4$, $60^\circ$ & $-0.08231846186746$ & $-0.08231846186746$\\
 & Hayward, $r=3.1$, $90^\circ$ & $-0.08168373732763$ & $-0.08168373732763$\\
 & KN $q=0.5$, $r=3.5$, $60^\circ$ & $-0.0790432393693279$ & $-0.0790432393693263$\\
 & KN $q=0.8$, $r=3.0$, $45^\circ$ & $-0.136651895747247$ & $-0.136651895747233$\\
on-axis $a^2$, $a^4$ coeff. & Bardeen, $r=3.4$ & $-0.0938935339679$, $-0.0161005095929$ & same to all digits\\
 & Bardeen, $r=3.1$ & $-0.114811334175$, $-0.0214112232194$ & same to all digits\\
$\langle L\rangle/(aE)$ & Bardeen, $r=3.4$, $60^\circ$ & $-1.01983100502$ & $-1.01983100502$\\
 & KN $q=0.5$, $r=3.5$, $60^\circ$ & $-0.920454545455$ & $-0.920454545455$\\
$a^3$ coeff.\ of $\langle L\rangle/E$, Kerr & $r=7/2$, $C=1/4$ & $-0.0825714285714285714$ & $-0.0825714285714285714$\\
 & $r=10/3$, $C=4/5$ & $-0.030955078125$ & $-0.030955078125$\\
 & $r=3.55$, $40^\circ$ & $-0.05225970864164288$ & $-0.05225970864164288$\\
$a^4$ coeff.\ of $\sigma/\sigma_0$, Kerr, eq.~\eqref{eq:kerr4} & $r=3.55$, $40^\circ$ & $-0.01290083149371615$ & $-0.01290083149371615$\\
 & $r=3.15$, $70^\circ$ & $-0.01428654979243309$ & $-0.01428654979243309$\\
\bottomrule
\end{tabular}}
\caption{Selected comparisons between the formulas and direct high-precision computations (values from
\texttt{verify\_paperA.log}, except the Hayward $50^\circ$ shadow row and the Bardeen $r=3.1$ and Hayward $r=3.4$ capture rows,
which come from the earlier scripts \texttt{univ\_check\_light.py} and \texttt{univ\_check\_mass.py}).
Bardeen: $g=0.4$; Hayward: $l=0.5$ ($M=1$). The direct computations use the exact contour or
the exact capture boundary, not the expansions.}
\label{tab:checks}
\end{table}

\section{Relation to previous work}\label{sec:prior}
To the best of our knowledge, based on searches of arXiv, INSPIRE and the wider web (October 2026), the
status of the results is as follows.

\emph{Previously known.}
(i) The exact shadow contour for metrics of the form \eqref{eq:metric}~\cite{Tsukamoto2018,Shaikh2019},
and the exact polar shadow radius for general integrable spinning spacetimes~\cite{Salehi2024}.
(ii) For Kerr, the order-$a^2$ shape of the shadow (an ellipse displaced by $2a\sin\vartheta_o$)~\cite{Bozza2006,GrallaLupsasca2020shape,PerlickTsupko2022},
hence the $a^2$ term of \eqref{eq:kerrA}; the edge-on area to order $a^2$ is stated explicitly in~\cite{Hassanabadi2026}.
(iii) Kobialko and Gal'tsov~\cite{KobialkoGaltsov2026} developed a perturbation theory for shadows of Kerr-like spacetimes,
expanding around Schwarzschild to fifth order in the spin and to second order in a small parameter measuring the deviation from Kerr.
Their results, through their areal definition of the mean radius, yield the order-$a^4$ Kerr area of \eqref{eq:kerrA} and the
Kerr--Newman double series to combined fourth order.
(iv) For Kerr, the order-$a^2$ capture cross-section at any speed and direction, and the swallowed angular momentum, follow
from the slow-rotation expansion of the critical Carter constant to $O(a^3)$ by Mach, Momennia and Sarbach~\cite{MachMomenniaSarbach2026},
who give the direction-averaged accretion rates explicitly (see also~\cite{MachMomenniaSarbach2026b}); Karydas, Vicente and Bertone
computed the direction- and speed-dependent coefficients through order $a^4$ as numerical fits~\cite{Karydas2026}. The
Schwarzschild cross-section $\sigma_0(v)$ is classical~\cite{Unruh1976}, and the Reissner--Nordstr\"om one at arbitrary speed was
studied in~\cite{Benone2014}.

\emph{New in this paper, as far as we could determine.}
(a) The formulas exact in $f$: Results~\ref{res:A2}--\ref{res:torque}, including the order-$a^4$ shadow formula, the
$a^6$ and $a^8$ coefficients, the capture cross-section of massive particles for general $f$, its on-axis continuation to
$a^4$ and the angular-momentum law $-a\sin^2\vartheta_o(1-f)/f$.
(b) For Kerr: the $a^6$ and $a^8$ shadow-area coefficients at general inclination, the full order-$a^4$ capture term
\eqref{eq:kerr4} (numerically identified), and the explicit direction-dependent order-$a^3$ swallowed angular momentum
\eqref{eq:kerrL3} (numerically identified, and derivable in principle from~\cite{MachMomenniaSarbach2026}).
(c) For Kerr--Newman: the closed forms exact in the charge, eqs.~\eqref{eq:KNsig}--\eqref{eq:KNA}.


\section{Discussion}
The formulas above reduce the small-spin behaviour of the shadow area, the capture cross-section
and the accreted angular momentum of any Kerr-like black hole of the form \eqref{eq:metric} to the
evaluation of $f$ and a few derivatives at one radius. Two limitations should be kept in mind.
First, the metrics \eqref{eq:metric} are generated by an algorithm and are in general not solutions of
any particular set of field equations; the results apply to them as geometries. Rotating metrics
of other forms (for example the Kerr--Sen metric, whose $\Sigma$ differs) are not covered. Second,
the results are expansions in $a$; their accuracy for rapidly rotating holes should be judged from
the size of the successive terms, which we provide up to $a^8$ for the shadow.

\section*{AI disclosure statement}
This work was carried out by the author with substantial assistance from an AI system, Claude
(Anthropic), used through the Claude Code software. Under the author's direction, the AI system
proposed the research directions, wrote and ran the computer-algebra and numerical code (Python
with SymPy and mpmath), performed the derivations and checks reported here, searched the
literature, and drafted the text of this paper. The author directed the work, set the requirements
(verification of every result, honest novelty assessment, no fabricated citations) and is
responsible for its publication. All results were checked numerically as described in
Section~\ref{sec:checks}; nevertheless, readers should treat the novelty assessment, which is based on
literature searches, with appropriate caution.

\begin{sloppypar}
\section*{Code and data availability}\label{sec:code}
All code used to derive and check the results, including the machine-readable $a^6$ and $a^8$
coefficients, is available at \url{https://github.com/JacobGoodchild/findformula2}
(directory \texttt{code/}; scripts \texttt{univ\_a2.py}, \texttt{univ\_a4\_fast.py},
\texttt{univ\_polar.py}, \texttt{univ\_check\_light.py}, \texttt{univ\_check\_mass.py},
\texttt{univ\_torque\_check.py}, \texttt{kn\_anyv\_a2.py}, \texttt{kn\_a4\_check.py}, and the
verification script \texttt{papers/paperA/verify\_paperA.py}).
\end{sloppypar}

\begin{thebibliography}{99}
\bibitem{EHT2019} Event Horizon Telescope Collaboration (K.~Akiyama et al.), ``First M87 Event Horizon Telescope Results. I. The Shadow of the Supermassive Black Hole'', Astrophys.\ J.\ Lett.\ \textbf{875}, L1 (2019), arXiv:1906.11238.
\bibitem{Bardeen1973} J.~M.~Bardeen, ``Timelike and null geodesics in the Kerr metric'', in \emph{Black Holes (Les Astres Occlus)}, Les Houches 1972, eds.\ C.~DeWitt and B.~S.~DeWitt (Gordon and Breach, New York, 1973), pp.\ 215--240.
\bibitem{Chandrasekhar1983} S.~Chandrasekhar, \emph{The Mathematical Theory of Black Holes} (Clarendon Press, Oxford, 1983).
\bibitem{NewmanJanis1965} E.~T.~Newman and A.~I.~Janis, ``Note on the Kerr spinning-particle metric'', J.\ Math.\ Phys.\ \textbf{6}, 915--917 (1965).
\bibitem{AzregAinou2014} M.~Azreg-A\"{\i}nou, ``Generating rotating regular black hole solutions without complexification'', Phys.\ Rev.\ D \textbf{90}, 064041 (2014), arXiv:1405.2569.
\bibitem{Tsukamoto2018} N.~Tsukamoto, ``Black hole shadow in an asymptotically-flat, stationary, and axisymmetric spacetime: The Kerr-Newman and rotating regular black holes'', Phys.\ Rev.\ D \textbf{97}, 064021 (2018), arXiv:1708.07427.
\bibitem{Shaikh2019} R.~Shaikh, ``Black hole shadow in a general rotating spacetime obtained through Newman-Janis algorithm'', Phys.\ Rev.\ D \textbf{100}, 024028 (2019), arXiv:1904.08322.
\bibitem{Cardoso2009} V.~Cardoso, A.~S.~Miranda, E.~Berti, H.~Witek and V.~T.~Zanchin, ``Geodesic stability, Lyapunov exponents and quasinormal modes'', Phys.\ Rev.\ D \textbf{79}, 064016 (2009), arXiv:0812.1806.
\bibitem{Bozza2006} V.~Bozza, F.~De~Luca and G.~Scarpetta, ``Kerr black hole lensing for generic observers in the strong deflection limit'', Phys.\ Rev.\ D \textbf{74}, 063001 (2006), arXiv:gr-qc/0604093.
\bibitem{GrallaLupsasca2020shape} S.~E.~Gralla and A.~Lupsasca, ``Observable shape of black hole photon rings'', Phys.\ Rev.\ D \textbf{102}, 124003 (2020), arXiv:2007.10336.
\bibitem{PerlickTsupko2022} V.~Perlick and O.~Yu.~Tsupko, ``Calculating black hole shadows: Review of analytical studies'', Phys.\ Rept.\ \textbf{947}, 1--39 (2022), arXiv:2105.07101.
\bibitem{KobialkoGaltsov2026} K.~Kobialko and D.~Gal'tsov, ``Perturbation theory for gravitational shadows in Kerr-like spacetimes'', Phys.\ Rev.\ D \textbf{113}, 104007 (2026), arXiv:2512.24126.
\bibitem{Hassanabadi2026} H.~Hassanabadi, M.~R.~R.~Good, S.~Zare, O.~Luongo and F.~Kafikang, ``Optical and thermodynamic properties of Kerr-Bertotti-Robinson black holes'', arXiv:2607.11979 (2026).
\bibitem{Salehi2024} K.~Salehi, A.~E.~Broderick and B.~Georgiev, ``Photon rings and shadow size for general axisymmetric and stationary integrable spacetimes'', Astrophys.\ J.\ \textbf{966}, 143 (2024), arXiv:2311.01495.
\bibitem{MachMomenniaSarbach2026} P.~Mach, M.~Momennia and O.~Sarbach, ``Accretion of a Vlasov gas by a Kerr black hole'', Phys.\ Rev.\ D \textbf{113}, 044068 (2026), arXiv:2508.04783.
\bibitem{MachMomenniaSarbach2026b} P.~Mach, M.~Momennia and O.~Sarbach, ``Bondi-type accretion onto a Kerr black hole in the kinetic regime'', Phys.\ Rev.\ D \textbf{113}, L041505 (2026), arXiv:2508.20189.
\bibitem{Karydas2026} T.~K.~Karydas, R.~Vicente and G.~Bertone, ``Mass and spin coevolution of black holes inspiralling through dark matter'', Phys.\ Rev.\ D \textbf{113}, 043057 (2026), arXiv:2510.13604.
\bibitem{Unruh1976} W.~G.~Unruh, ``Absorption cross-section of small black holes'', Phys.\ Rev.\ D \textbf{14}, 3251 (1976).
\bibitem{Benone2014} C.~L.~Benone, E.~S.~de~Oliveira, S.~R.~Dolan and L.~C.~B.~Crispino, ``Absorption of a massive scalar field by a charged black hole'', Phys.\ Rev.\ D \textbf{89}, 104053 (2014), arXiv:1404.0687.
\bibitem{AyonBeatoGarcia2000} E.~Ay\'on-Beato and A.~Garc\'{\i}a, ``The Bardeen model as a nonlinear magnetic monopole'', Phys.\ Lett.\ B \textbf{493}, 149 (2000), arXiv:gr-qc/0009077.
\bibitem{Hayward2006} S.~A.~Hayward, ``Formation and evaporation of regular black holes'', Phys.\ Rev.\ Lett.\ \textbf{96}, 031103 (2006), arXiv:gr-qc/0506126.

\end{thebibliography}
\end{document}
